\documentclass[11pt,a4paper]{article}
\usepackage[margin=25mm]{geometry}
\usepackage{amsmath,amssymb,amsthm}
\newtheorem{lemma}{Lemma}\newtheorem{proposition}[lemma]{Proposition}
\title{Implementing a central twist on two permutation domains}
\author{Human Compression project: F006 continuation}\date{}
\begin{document}\maketitle
The complete proof of \emph{Enumeration by stabilizers of initial strings}
remains a prerequisite. Its orbit recursion, intersection test and necessary
completion test are unchanged. This note replaces a large permutation
representation of the automorphism group by a smaller conjugation action.
It makes no claim that the exceptional census is complete.

\begin{lemma}
Let $G\leq\operatorname{Sym}(\Omega)$ be finite. Suppose $z\in Z(G)$
has order two and $\chi:G\to\{0,1\}$ is an additive character with
$\chi(z)=0$. Then $\phi(g)=z^{\chi(g)}g$ is an involutory automorphism.
On $\Omega\times\{0,1\}$ define
\[
 \iota(g)(x,0)=(g(x),0),\qquad
 \iota(g)(x,1)=(\phi(g)(x),1),\qquad s(x,i)=(x,1-i).
\]
The map $\iota$ is an injective homomorphism, and conjugation by $s$
implements $\phi$ on $\iota(G)$. Put
$A=\langle\iota(G),s\rangle$. The image of the conjugation action of
$A$ on $\iota(G)$ is exactly
$\langle\operatorname{Inn}(G),\phi\rangle$.
\end{lemma}
\begin{proof}
Centrality and additivity give
$\phi(gh)=\phi(g)\phi(h)$. Since $\chi(z)=0$, applying $\phi$ twice
returns $g$. Thus it is an automorphism, including when it is the identity.
The two coordinate actions define a homomorphism; its first coordinate is
the given faithful action of $G$, so the homomorphism is injective.
Swapping the two domains interchanges the actions $g$ and $\phi(g)$.
Because $\phi^2=1$, this gives
$s^{-1}\iota(g)s=\iota(\phi(g))$. Consequently $\iota(G)$ is normal
in $A$. Conjugation by its elements gives the inner automorphisms, and
conjugation by the additional generator gives $\phi$; these generate
exactly the stated image.
\end{proof}

\begin{proposition}
Suppose exact group operations establish
\[
 |A:C_A(\iota(G))|=|\operatorname{Aut}(G)|.                 \tag{1}
\]
In the retained enumeration algorithm, one may use $A$ acting by
conjugation on the involutions of $\iota(G)$. For an already chosen
tuple $t$, use $C_A(t)$ in place of its automorphism stabilizer.
The resulting search records precisely the same abstract generating
string C-markings, up to isomorphism, as the full automorphism action.
The same statement applies to its reversal and connected-string counts.
\end{proposition}
\begin{proof}
Normality makes the conjugation image a subgroup of the full automorphism
group. Its kernel is exactly $C_A(\iota(G))$, so (1) makes the image the
whole finite automorphism group. The map $A\to\operatorname{Aut}(G)$
is therefore surjective. For each tuple $t$, the inverse image of its
pointwise automorphism stabilizer is exactly $C_A(t)$: membership means
fixing every chosen group element under conjugation. Hence both groups
have the same orbits on all candidate involutions. This remains true at
every successive prefix. The retained completeness proof applies without
change. The nontrivial kernel acts trivially throughout; it neither creates
nor removes a candidate orbit. Surjectivity also gives the same existence
test for an automorphism carrying a full tuple to its reversal. Connectedness
depends only on the orders of adjacent products and is preserved by $\iota$.
\end{proof}

\paragraph{The actual finite guards.}
The implementation takes $\chi$ to be permutation parity. It checks that
the original model has its stated order, that its centre has order two,
and that the unique nonidentity central element is even. It builds the
explicit permutations above, verifies their generated order, normality,
and conjugation identity on every original generator, and checks (1)
against the exact order returned by GAP's full automorphism computation.
It then uses the unchanged conjugation recursion and completion pruning.
A failure of any guard stops the search; no assertion that this twist
always supplies every outer automorphism is used.

For the retained $E_8$ reflection model the guarded calculation gives
original degree $240$, new degree $480$, group order $696729600$,
overgroup order $1393459200$, conjugation kernel order $2$ and full
automorphism order $696729600$. This avoids constructing its automorphism
permutations on all $199951$ involutions. Those involutions still form the
complete candidate set. The original model constructor, exact group
operations and complete census output remain required computational
premises. The group orders in this paragraph are checked finite facts,
not replacements for the guards or for the still unfinished census.

As a small control, take
$G=\langle(1\ 2),(1\ 2\ 3),(4\ 5)(6\ 7)\rangle$.
The same construction gives degree fourteen and a full conjugation action
of order twelve. Both algorithms enumerate its entire seven-involution
set and the same complete generating markings: one ordered class of rank
two and two of rank three. This is a check on the implementation; the
uniform justification is the two proofs above.

\paragraph{Attribution and limits.}
The central automorphism and induced permutation action are standard group
constructions; no priority is claimed for them. Their explicit guarded use
here is an implementation improvement for the retained stabilizer search,
whose Hartley--Hulpke and Leemans--Mixer provenance remains unchanged.
A complete exceptional-group computation and distinct mathematical review
are still required. This note supplies no whole-original, open-question,
Lean or compression-factor credit.
\end{document}
